\documentclass[leqno]{article}
% Shared QBT/Soln pack preamble for resources.danieljsmith.org.
\usepackage{graphicx}
\usepackage{dsfont}
\usepackage{mathtools}
\usepackage{amssymb}
\usepackage{amsmath}
\usepackage{mleftright}
\usepackage{amsthm}
\usepackage{float}
\usepackage{ragged2e}
\usepackage{tensor}
\usepackage{fancyhdr}
\usepackage{empheq}
\usepackage[most]{tcolorbox}
\usepackage{enumitem}
\usepackage{dirtytalk}
\usepackage{marginnote}
\usepackage{microtype}
\usepackage[
  a4paper,
  left=0.8in,
  right=1in,
  top=1in,
  bottom=0.5in,
  headheight=14pt,
  headsep=0.4in
]{geometry}
\setlength{\parindent}{0cm}

% TikZ / PGFPlots
\usepackage{pgfplots}
\pgfplotsset{compat=1.18}
\usepgfplotslibrary{fillbetween, polar}
\usetikzlibrary{calc, positioning}
\usetikzlibrary{angles, quotes}
\usetikzlibrary{arrows.meta}
\usetikzlibrary{decorations.pathreplacing, decorations.markings}
\usetikzlibrary{patterns}
\usetikzlibrary{intersections}

% Links and contents. The contents follow the TMUA papers: compact entries,
% no dot leaders, and each question's mark total in grey after its label.
\usepackage{tocloft}
\tocloftpagestyle{fancy}
\setlength{\cftbeforesubsecskip}{2pt}
\setlength{\cftsubsecindent}{1.5em}
\setlength{\cftsubsecnumwidth}{0pt}
\renewcommand{\cftsubsecfont}{\small\color{black}}
\renewcommand{\cftsubsecpagefont}{\small\color{black}}
\renewcommand{\cftsubsecleader}{\hfill}
\usepackage[colorlinks=true,linkcolor=black,citecolor=magenta,urlcolor=black]{hyperref}%
\urlstyle{same}

% Pack macros
% \marks also adds its value to the current question's total in the .aux
% file, so the contents show each total from the second pass. The write
% happens at shipout, so a box typeset more than once counts once. The
% pack's mark total for the contents-page title is summed the same way.
\makeatletter
\newcount\djs@totalmarks
\newcommand{\djsAddMarks}[2]{%
  \global\advance\djs@totalmarks#2\relax
  \@ifundefined{djs@marks@#1}%
    {\expandafter\xdef\csname djs@marks@#1\endcsname{#2}}%
    {\expandafter\xdef\csname djs@marks@#1\endcsname
      {\the\numexpr\csname djs@marks@#1\endcsname+#2\relax}}}
\RenewDocumentCommand{\marks}{o m}{%
  \protected@write\@auxout{}{\string\djsAddMarks{\number\value{enumi}}{#2}}%
  \IfValueTF{#1}%
    {\marginnote{\raggedright\textbf{[#2]}}[#1]}%
    {\marginnote{\raggedright\textbf{[#2]}}}%
}
\makeatother

\newcommand{\vspacerule}[2]{%
  \par\vspace*{#1}%
  \noindent\hrulefill\par
  \vspace*{#2}%
}

\definecolor{scarlet}{RGB}{205, 40, 75}
% Site colours for the contents-page title: the accent, the rule colour, a
% stronger rule colour for the path separator, and the muted text colour.
\definecolor{djsBurgundy}{RGB}{118, 59, 67}
\definecolor{djsHairline}{RGB}{216, 213, 205}
\definecolor{djsSlash}{RGB}{170, 166, 158}
\definecolor{djsMuted}{RGB}{96, 101, 106}

% Worked-solution box, adapted from the TMUA worked-solutions box. A white breakable box with a left scarlet rule and no surrounding
% frame or tint. A part that breaks fills the rest of its page, so the rule
% runs to the bottom margin; the unbroken or last part ends in a short
% horizontal foot rule so a solution visibly stops. Unlike the TMUA box there
% is no answer label. The box does not grow to the left, so the rule sits at
% the text edge; growing to the right keeps the usable line width close to the
% old tcolorbox Solution box.
% Spacing inside the box: paragraphs are 6pt apart, so source blank lines show
% (solutions put one after each display and none between a lead-in and its
% display). Every display gets 6pt above and 10pt below, whatever the length
% of its lead-in line; align rows are spread by an 8pt \jot; solution parts
% are 14pt apart and sub-parts 8pt; \left and \right space like ordinary
% brackets (mleftright); no single line is stranded at a page break; and
% neither the Solution label nor the line before a list ends a page. No
% skip can shrink, so full pages are not squeezed.
\newtcolorbox{djsSolution}{%
  enhanced, breakable, frame hidden, sharp corners, colback=white,
  height fixed for=first and middle, height fill,
  extras unbroken and last={natural height},
  borderline west={2.2pt}{0pt}{scarlet},
  left=11pt, right=0pt, top=5pt, bottom=13pt,
  before skip=13pt, after skip=4pt,
  grow to right by=10mm,
  overlay unbroken and last={%
    \fill[scarlet] (frame.south west)
      rectangle ([xshift=3.5cm, yshift=2.2pt]frame.south west);},
  before upper={{\color{scarlet}\bfseries Solution}\par\nobreak\vspace{4pt}%
    \setlength{\parskip}{6pt plus 2pt}%
    \setlength{\abovedisplayskip}{6pt plus 2pt}%
    \setlength{\abovedisplayshortskip}{6pt plus 2pt}%
    \setlength{\belowdisplayskip}{10pt plus 2pt}%
    \setlength{\belowdisplayshortskip}{10pt plus 2pt}%
    \setlength{\jot}{8pt}%
    \setlist{parsep=6pt plus 2pt, topsep=6pt plus 2pt, beginpenalty=10000}%
    \setlist[enumerate,1]{itemsep=14pt plus 2pt}%
    \setlist[enumerate,2]{itemsep=8pt plus 2pt}%
    \mleftright
    \clubpenalty=10000 \widowpenalty=10000 \displaywidowpenalty=10000\relax}}

% Optional smaller figure inside a solution box. \scalebox shrinks node text
% with the drawing. Not applied to existing figures.
\newcommand{\SolnFigureScale}{0.85}
\newsavebox{\djsSolnFigureBox}
\newenvironment{djsSolutionFigure}
  {\par\vspace{4pt}\begin{center}\begin{lrbox}{\djsSolnFigureBox}}
  {\end{lrbox}\scalebox{\SolnFigureScale}{\usebox{\djsSolnFigureBox}}%
   \end{center}\vspace{2pt}\par}

\newcommand{\cxl}[2]{%
  \tikz[baseline=(X.base)]{%
    \node[inner sep=0pt, outer sep=0pt, text=#1] (X) {$#2$};
    \draw[#1, line width=0.4pt] (X.south west)--(X.north east);
  }%
}

% Page-style hook run at the contents page and each question's opening page.
% A no-op for QBT packs; \djsSolnHeader points it at the fancy header.
\newcommand{\djsQuestionPageStyle}{}

% A contents entry: "Question N" in a box wide enough for two digits, so the
% mark totals line up, then the total in grey. Bookmarks keep "Question N".
\makeatletter
\newlength{\djsTocLabelWidth}
\newcommand{\djsTocMarks}[1]{%
  \@ifundefined{djs@marks@#1}{}{%
    {\footnotesize\itshape\color{black!45}%
     \csname djs@marks@#1\endcsname~mark%
     \ifnum\csname djs@marks@#1\endcsname=1 \else s\fi}}}
\newcommand{\djsTocEntry}[1]{%
  \texorpdfstring
    {\settowidth{\djsTocLabelWidth}{Question 00}%
     \makebox[\dimexpr\djsTocLabelWidth+0.9em\relax][l]{Question #1}%
     \djsTocMarks{#1}}%
    {Question #1}}
\makeatother

% The question count for the contents-page title is written to the .aux at
% the end of the document, so it also appears from the second pass.
\newcounter{djsQuestions}
\newcommand{\questionitem}{%
  \djsQuestionPageStyle
  \item\phantomsection
  \stepcounter{djsQuestions}%
  \addcontentsline{toc}{subsection}{\protect\djsTocEntry{\number\value{enumi}}}%
}
\makeatletter
\newcommand{\djsSetQuestionCount}[1]{\gdef\djs@questioncount{#1}}
\AtEndDocument{\immediate\write\@auxout
  {\string\djsSetQuestionCount{\number\value{djsQuestions}}}}
\makeatother

\renewcommand{\contentsname}{Questions}

% \djsPackHeader[<link content tag>]{<left header>}; the tag says which
% document a click on the header link came from.
\newcommand{\djsPackHeader}[2][questions]{%
  \pagestyle{fancy}%
  \fancyhead{}%
  \fancyfoot{}%
  \renewcommand{\headrulewidth}{0.9pt}%
  \lhead{#2}%
  \chead{}%
  \rhead{\href{https://danieljsmith.org/?utm_source=pdf&utm_campaign=further-maths&utm_content=#1}{danieljsmith.org}}%
}

% The header's topic is also the contents-page title, under which a label
% says which booklet this is.
\newcommand{\djsPackTitle}{}
\newcommand{\djsBookletLabel}{}
\newcommand{\djsQbtHeader}[1]{%
  \djsPackHeader{#1}%
  \renewcommand{\djsPackTitle}{#1}%
  \renewcommand{\djsBookletLabel}{Questions by topic}%
}

% The colophon at the foot of page one, independent of the page style: the
% contents-page title's hairline led by a short scarlet rule, then the date on
% the left and a link to the site on the right, all inside the bottom margin.
% Its sizes are fixed, so it matches the TMUA papers' colophon. The document
% supplies the date; recompiling never changes it.
\makeatletter
\newcommand{\djsLastUpdated}[1]{%
  \AtBeginDocument{\AddToHookNext{shipout/foreground}{\djs@colophon{#1}}}}
\newcommand{\djs@colophon}[1]{%
  \put(0.8in,-\dimexpr\paperheight-0.47in\relax){%
    \makebox[0pt][l]{\color{djsHairline}%
      \rule[0.9pt]{\dimexpr\paperwidth-1.8in\relax}{0.6pt}}%
    \makebox[0pt][l]{\color{scarlet}\rule{3.5cm}{2.2pt}}}%
  \put(0.8in,-\dimexpr\paperheight-0.32in\relax){\makebox[0pt][l]{%
    \fontsize{9}{11}\selectfont Last updated #1}}%
  \put(\dimexpr\paperwidth-1in\relax,-\dimexpr\paperheight-0.32in\relax){%
    \makebox[0pt][r]{\fontsize{7}{8.5}\selectfont\color{djsMuted}%
      Latest version at \href{https://resources.danieljsmith.org/}%
        {\color{black}resources.danieljsmith.org}}}}
\makeatother

% Solutions packs show the header only on the contents page and each
% question's opening page, so it never appears part-way through a solution.
% \thispagestyle marks the next page shipped out, which is the question's
% page because every solution question starts after a \newpage.
\newcommand{\djsSolnHeader}[1]{%
  \djsPackHeader[solutions]{\textit{Solutions} - #1}%
  \pagestyle{empty}%
  \renewcommand{\djsQuestionPageStyle}{\thispagestyle{fancy}}%
  \renewcommand{\djsPackTitle}{#1}%
  \renewcommand{\djsBookletLabel}{Worked solutions}%
}

% \djsSitePath{<strand>}{<topic>}, called once after the header, names where
% the pack sits on the site; the contents-page title prints it above the pack
% title, e.g. CORE PURE / FURTHER CALCULUS. Without it the line is left out.
\makeatletter
\newif\ifdjs@sitepath
\newcommand{\djs@strand}{}
\newcommand{\djs@section}{}
\newcommand{\djsSitePath}[2]{%
  \djs@sitepathtrue
  \renewcommand{\djs@strand}{#1}%
  \renewcommand{\djs@section}{#2}%
}
\makeatother

% Further Mechanics packs call \djsMechanicsPack in their preamble, so that
% \djsFrontMatter prints the default value of g above the contents. The box
% keeps its own space above it (before skip is the whole gap below the title,
% not an addition to it). Without the box, \djsGravityNote is just extra space
% before the contents heading, bringing the gap above it close to the gap
% between the box and the heading on Further Mechanics pages.
\newcommand{\djsGravityNote}{\vspace{14pt}}
\newcommand{\djsMechanicsPack}{%
  \renewcommand{\djsGravityNote}{%
    \begin{tcolorbox}[sharp corners, colback=white, colframe=black,
      boxrule=1pt, left=8pt, right=8pt, top=6pt, bottom=6pt,
      before skip=20pt]
    {\bfseries Instructions}\par\vspace{4pt}
    Unless a question states otherwise, take $g = 9.8\,\mathrm{m\,s^{-2}}$
    wherever a numerical value of $g$ is required. Solutions also give final
    answers for $g = 10\,\mathrm{m\,s^{-2}}$ where these differ.
    \end{tcolorbox}%
    \vspace{16pt}}%
}

% Contents-page title: the site path in spaced burgundy capitals, the pack
% title, a hairline led by a short scarlet rule (the solution box's foot rule),
% then the booklet label with the question and mark counts. The counts come
% from the .aux, so they appear from the second pass.
\makeatletter
\newsavebox{\djs@eyebrowbox}
\newcommand{\djs@caps}[1]{\textls[180]{\MakeUppercase{#1}}}
\newcommand{\djs@eyebrow}{%
  \color{djsBurgundy}\footnotesize\bfseries
  \djs@caps{\djs@strand}%
  {\color{djsSlash}\hspace{0.55em}/\hspace{0.55em}}%
  \djs@caps{\djs@section}}
\newcommand{\djs@dot}{\hspace{0.45em}\textperiodcentered\hspace{0.45em}}
\newcommand{\djs@counts}{%
  \@ifundefined{djs@questioncount}{}{%
    \djs@questioncount~question\ifnum\djs@questioncount=1 \else s\fi
    \ifnum\djs@totalmarks>0
      \djs@dot\the\djs@totalmarks~mark\ifnum\djs@totalmarks=1 \else s\fi
    \fi}}
\newcommand{\djsContentsTitle}{%
  \ifdjs@sitepath
    \sbox\djs@eyebrowbox{\djs@eyebrow}%
    \ifdim\wd\djs@eyebrowbox>\linewidth
      \@latex@warning@no@line{The \string\djsSitePath\space line is wider
        than the text, so it wraps. Shorten its topic}%
    \fi
    {\raggedright\djs@eyebrow\par}%
    \vspace{7pt}%
  \fi
  {\raggedright\fontsize{24.88}{28}\selectfont\djsPackTitle\par}%
  \vspace{9pt}%
  \noindent\makebox[0pt][l]{\color{djsHairline}\rule[0.9pt]{\linewidth}{0.6pt}}%
  {\color{scarlet}\rule{3.5cm}{2.2pt}}\par
  \vspace{5pt}%
  {\small\color{djsMuted}\djsBookletLabel\hfill\djs@counts\par}%
  \vspace{8pt}%
}
\makeatother

\newcommand{\djsFrontMatter}{%
  \djsQuestionPageStyle
  \vspace*{6pt}%
  \djsContentsTitle
  \djsGravityNote
  \tableofcontents
  \newpage
}


\djsSolnHeader{Dimensional Analysis}
\djsSitePath{Further Mechanics}{Miscellaneous Mechanics}
\djsMechanicsPack

\djsLastUpdated{27 October 2026}
\begin{document}
\djsFrontMatter
\begin{enumerate}[
  label=\textbf{\arabic*.},
  leftmargin=*,
  topsep=0pt,
  partopsep=0pt,
  parsep=0pt
]

% ---- Question 1 ----
\questionitem
% [Source: _QBT___Solns__Dimensional_Analysis, Q1]

\begin{enumerate}[label=\textbf{(\alph*)}]
  \item State the dimensions of force. \marks{1} \\
  \item The braking distance of a vehicle can be modelled by the formula
  \[
  d = m^{\alpha} u^{\beta} R^{\gamma}
  \] 
  where \\

  $m$ is the mass of the vehicle in $\mathrm{kg}$

  $u$ is the initial speed in $\mathrm{m\,s^{-1}}$

  $R$ is the constant braking force in $\mathrm{N}$ \\
  
  Use dimensional analysis to find the values of $\alpha$, $\beta$ and $\gamma$. \marks{4} \\
\end{enumerate}

\vspace*{10pt}

\begin{djsSolution}
\begin{enumerate}[label=\textbf{(\alph*)}, start=1]
  \item Force is mass $\times$ acceleration, so
  \[
  [F]=\mathrm{M}\times\mathrm{LT^{-2}}=\mathrm{MLT^{-2}}
  \]

  \item Write the dimensions of each quantity:
  \[
  [d]=\mathrm{L},\qquad [m]=\mathrm{M},\qquad [u]=\mathrm{LT^{-1}},\qquad [R]=\mathrm{MLT^{-2}}
  \]

  Given
  \[
  d=m^{\alpha} u^{\beta} R^{\gamma}
  \]

  the dimensions on the right-hand side are
  \begin{align*}
  [m^{\alpha} u^{\beta} R^{\gamma}]
  &=\mathrm{M}^{\alpha}(\mathrm{LT^{-1}})^{\beta}(\mathrm{MLT^{-2}})^{\gamma}\\
  &=\mathrm{M}^{\alpha}\mathrm{L}^{\beta}\mathrm{T}^{-\beta}\times\mathrm{M}^{\gamma}\mathrm{L}^{\gamma}\mathrm{T}^{-2\gamma}\\
  &=\mathrm{M}^{\alpha+\gamma}\mathrm{L}^{\beta+\gamma}\mathrm{T}^{-\beta-2\gamma}
  \end{align*}

  This must equal the dimensions of $d$, so
  \[
  \mathrm{M}^{\alpha+\gamma}\mathrm{L}^{\beta+\gamma}\mathrm{T}^{-\beta-2\gamma}=\mathrm{L}=\mathrm{M^0L^1T^0}
  \]

  Equating powers of $\mathrm{M}$, $\mathrm{L}$ and $\mathrm{T}$ gives
  \[
  \alpha+\gamma=0,\qquad \beta+\gamma=1,\qquad -\beta-2\gamma=0
  \]

  The equation $-\beta-2\gamma=0$ gives
  \[
  \beta=-2\gamma
  \]

  Substitute into $\beta+\gamma=1$:
  \begin{align*}
  -2\gamma+\gamma&=1\\
  -\gamma&=1\\
  \gamma&=-1
  \end{align*}

  Then
  \[
  \beta=-2(-1)=2
  \]

  and $\alpha+\gamma=0$ gives
  \[
  \alpha+(-1)=0
  \]

  so
  \[
  \alpha=1
  \]

  Hence
  \[
  \alpha=1,\qquad \beta=2,\qquad \gamma=-1
  \]
\end{enumerate}
\end{djsSolution}


\newpage

% ---- Question 2 ----
\questionitem
% [Source: _QBT___Solns__Dimensional_Analysis, Q2]

A small satellite moves in a circular orbit of radius $r$ about the centre of a much more massive planet of mass $M$. It is suggested that the period $T$ of the orbit depends only on $G$, $M$ and $r$, and that
\[
T = kG^{\alpha} M^{\beta} r^{\gamma}
\]
where $k$ is a dimensionless constant. \\

It is given that
\[
G = 6.67 \times 10^{-11}\,\mathrm{m}^3\,\mathrm{kg}^{-1}\,\mathrm{s}^{-2}
\]
\begin{enumerate}[label=\textbf{(\alph*)}, start=1]
  \item Write down the dimensions of $G$. \marks{1} \\
  \item By using dimensional analysis, determine the values of $\alpha$, $\beta$ and $\gamma$. \marks{4} \\
\end{enumerate}

\vspace*{10pt}

\begin{djsSolution}
\begin{enumerate}[label=\textbf{(\alph*)}, start=1]
  \item From
  \[
  G=6.67\times 10^{-11}\,\mathrm{m}^3\,\mathrm{kg}^{-1}\,\mathrm{s}^{-2}
  \]

  the dimensions of \(G\) are
  \[
  [G]=\mathrm{M^{-1}L^3T^{-2}}
  \]

  \item The constant \(k\) is dimensionless, so
  \[
  [T]=[G]^\alpha[M]^\beta[r]^\gamma
  \]

  Now
  \[
  [T]=\mathrm{T},\qquad [G]=\mathrm{M^{-1}L^3T^{-2}},\qquad [M]=\mathrm{M},\qquad [r]=\mathrm{L}
  \]

  so
  \begin{align*}
  \mathrm{T}&=(\mathrm{M^{-1}L^3T^{-2}})^\alpha\mathrm{M}^\beta\mathrm{L}^\gamma \\
   &=\mathrm{M}^{-\alpha+\beta}\mathrm{L}^{3\alpha+\gamma}\mathrm{T}^{-2\alpha}
  \end{align*}

  Equating powers of \(\mathrm{M}\), \(\mathrm{L}\) and \(\mathrm{T}\) with \(\mathrm{M^0L^0T^1}\):
  \begin{align*}
  -\alpha+\beta&=0 \\
  3\alpha+\gamma&=0 \\
  -2\alpha&=1
  \end{align*}

  The equation \(-2\alpha=1\) gives
  \[
  \alpha=-\frac12
  \]

  Then
  \[
  \beta=\alpha=-\frac12
  \]

  and
  \[
  \gamma=-3\alpha=\frac32
  \]

  Hence
  \[
  \alpha=-\frac12,\qquad \beta=-\frac12,\qquad \gamma=\frac32
  \]
\end{enumerate}
\end{djsSolution}


\newpage

% ---- Question 3 ----
\questionitem
% [Source: _QBT___Solns__Dimensional_Analysis, Q3]

A small probe is launched vertically upwards from the surface of a spherical moon. The radius of the moon is $R$ metres, and the probe is at a distance $r$ metres from the centre of the moon. The gravitational field strength $g$ at distance $r$ is modelled by
\[
g = \frac{\Gamma}{r^2}
\]
where $\Gamma$ is a constant associated with the moon. \\

The probe has initial speed $u$ at the surface and speed $v$ when it reaches distance $r$. These speeds are related by
\[
v^2 = u^2 - 2\Gamma\left(\frac{1}{R}-\frac{1}{r}\right)
\]
The work done against gravity in moving a probe of mass $m$ from the surface to distance $r$ is
\[
W = m\Gamma\left(\frac{1}{R}-\frac{1}{r}\right)
\]

\begin{enumerate}[label=\textbf{(\alph*)}]
  \item Find the dimensions of $\Gamma$. \marks{2}
  \item Use dimensional analysis to show that the equation for $v^2$ is dimensionally consistent. \marks{2}
  \item Find the dimensions and SI units of $W$. \marks{2}
\end{enumerate}

\vspace*{20pt}

\begin{djsSolution}
\begin{enumerate}[label=\textbf{(\alph*)}, start=1]
  \item From
  \[
  g=\frac{\Gamma}{r^2}
  \]

  and since
  \[
  [g]=\mathrm{LT^{-2}},\qquad [r]=\mathrm{L}
  \]

  we have
  \[
  [\Gamma]=[g][r^2]=\mathrm{LT^{-2}}\times\mathrm{L^2}=\mathrm{L^3T^{-2}}
  \]

  \item In
  \[
  v^2=u^2-2\Gamma\left(\frac{1}{R}-\frac{1}{r}\right)
  \]

  the dimensions of \(v^2\) and \(u^2\) are
  \[
  [v^2]=[u^2]=(\mathrm{LT^{-1}})^2=\mathrm{L^2T^{-2}}
  \]

  Also,
  \[
  \left[\frac{1}{R}\right]=\left[\frac{1}{r}\right]=\mathrm{L^{-1}}
  \]

  so
  \[
  \left[\left(\frac{1}{R}-\frac{1}{r}\right)\right]=\mathrm{L^{-1}}
  \]

  The factor \(2\) is dimensionless, so
  \[
  \left[2\Gamma\left(\frac{1}{R}-\frac{1}{r}\right)\right]=[\Gamma]\mathrm{L^{-1}}=\mathrm{L^3T^{-2}}\times\mathrm{L^{-1}}=\mathrm{L^2T^{-2}}
  \]

  Therefore every term in the equation has dimensions \(\mathrm{L^2T^{-2}}\), so the equation is dimensionally consistent.

  \item Using
  \[
  W=m\Gamma\left(\frac{1}{R}-\frac{1}{r}\right)
  \]

  with
  \[
  [m]=\mathrm{M},\qquad [\Gamma]=\mathrm{L^3T^{-2}},\qquad \left[\left(\frac{1}{R}-\frac{1}{r}\right)\right]=\mathrm{L^{-1}}
  \]

  gives
  \[
  [W]=\mathrm{M}\times\mathrm{L^3T^{-2}}\times\mathrm{L^{-1}}=\mathrm{ML^2T^{-2}}
  \]

  Its SI units are
  \[
  \mathrm{kg\,m^2\,s^{-2}}
  \]

  which is the joule, \(\mathrm{J}\).
\end{enumerate}
\end{djsSolution}


\newpage

% ---- Question 4 ----
\questionitem
% [Source: _QBT___Solns__Dimensional_Analysis, Q4]

The increase in pressure $p$ pascals between the surface and the bottom of a liquid of depth $h$ metres and density $\rho\,\mathrm{kg\,m^{-3}}$ is modelled by
\[
p = k\rho^{\alpha}h^{\beta}g
\]
where $g\,\mathrm{m\,s^{-2}}$ is the acceleration due to gravity and $k$ is a dimensionless constant. \\
\begin{enumerate}[label=\textbf{(\alph*)}]
  \item State the dimensions of $p$. \marks{1} \\
  \item Use dimensional analysis to find the value of $\alpha$ and the value of $\beta$. \marks{3} \\
\end{enumerate}

\vspace*{10pt}

\begin{djsSolution}
\begin{enumerate}[label=\textbf{(\alph*)}, start=1]
  \item Pressure is force per unit area.

  With \([\text{force}] = \mathrm{MLT^{-2}}\) and \([\text{area}] = \mathrm{L^2}\), this gives
  \[
  [p] = \frac{\mathrm{MLT^{-2}}}{\mathrm{L^2}} = \mathrm{ML^{-1}T^{-2}}
  \]

  \item Using
  \[
  p = k\rho^{\alpha}h^{\beta}g
  \]

  and noting that \(k\) is dimensionless gives
  \[
  [p] = [\rho]^{\alpha}[h]^{\beta}[g]
  \]

  Now
  \[
  [\rho] = \mathrm{ML^{-3}}, \qquad [h] = \mathrm{L}, \qquad [g] = \mathrm{LT^{-2}}
  \]

  So the dimensions of the right-hand side are
  \begin{align*}
  (\mathrm{ML^{-3}})^{\alpha}(\mathrm{L})^{\beta}(\mathrm{LT^{-2}})
  &= \mathrm{M}^{\alpha}\mathrm{L}^{-3\alpha}\mathrm{L}^{\beta}\mathrm{LT^{-2}} \\
  &= \mathrm{M}^{\alpha}\mathrm{L}^{-3\alpha+\beta+1}\mathrm{T^{-2}}
  \end{align*}

  But part (a) gives
  \[
  [p] = \mathrm{ML^{-1}T^{-2}}
  \]

  Equating powers of \(\mathrm{M}\), \(\mathrm{L}\) and \(\mathrm{T}\):
  \[
  \alpha = 1
  \]

  and
  \[
  -3\alpha + \beta + 1 = -1
  \]

  Substituting \(\alpha=1\):
  \begin{align*}
  -3(1) + \beta + 1 &= -1 \\
  \beta - 2 &= -1 \\
  \beta &= 1
  \end{align*}

  Hence
  \[
  \alpha = 1, \qquad \beta = 1
  \]
\end{enumerate}
\end{djsSolution}


\newpage

% ---- Question 5 ----
\questionitem
% [Source: _QBT___Solns__Dimensional_Analysis, Q5]

The force, $F$, exerted on a flat plate by a water jet is thought to depend on the density of the water, $\rho$, the cross-sectional area of the jet, $A$, and the speed of the water, $v$. You are given that $\rho$ is measured in $\mathrm{kg\,m^{-3}}$. \\
\begin{enumerate}[label=\textbf{(\alph*)}]
  \item Find the dimensions of density $\rho$. \marks{1} \\
\end{enumerate}
A model of the force suggests the relationship
\[
F = \frac12 \rho^\alpha A^\beta v^\gamma
\]
\begin{enumerate}[label=\textbf{(\alph*)}, resume]
  \item Explain whether or not it is possible to use dimensional analysis to verify that the constant $\frac12$ is correct. \marks{1} \\
  \item Use dimensional analysis to find the values of $\alpha$, $\beta$ and $\gamma$. \marks{4} \\
\end{enumerate}

\vspace*{10pt}

\begin{djsSolution}
\begin{enumerate}[label=\textbf{(\alph*)}, start=1]
  \item Density is mass per unit volume, so
\[
[\rho]=\frac{\mathrm{M}}{\mathrm{L^3}}=\mathrm{ML^{-3}}
\]


  \item No, it is not possible.

Dimensional analysis only checks that both sides of an equation have the same dimensions. The number \(\tfrac12\) is a dimensionless constant, so dimensional analysis cannot show whether \(\tfrac12\) is the correct numerical value.

  \item We use
\[
F=\frac12\rho^{\alpha}A^{\beta}v^{\gamma}
\]

The factor \(\tfrac12\) is dimensionless, so it does not affect the dimensions.

Now write the dimensions of each quantity:
\[
[F]=\mathrm{MLT^{-2}},\qquad [\rho]=\mathrm{ML^{-3}},\qquad [A]=\mathrm{L^2},\qquad [v]=\mathrm{LT^{-1}}
\]

So
\[
\mathrm{MLT^{-2}}=(\mathrm{ML^{-3}})^{\alpha}(\mathrm{L^2})^{\beta}(\mathrm{LT^{-1}})^{\gamma}
\]

Combining powers gives
\[
\mathrm{MLT^{-2}}=\mathrm{M}^{\alpha}\mathrm{L}^{-3\alpha+2\beta+\gamma}\mathrm{T}^{-\gamma}
\]

Equating powers of \(\mathrm{M}\), \(\mathrm{L}\) and \(\mathrm{T}\):
\begin{align*}
\text{For }\mathrm{M}:&\qquad \alpha=1 \\
\text{For }\mathrm{T}:&\qquad -\gamma=-2\quad\Longrightarrow\quad\gamma=2 \\
\text{For }\mathrm{L}:&\qquad -3\alpha+2\beta+\gamma=1
\end{align*}

Substitute \(\alpha=1\) and \(\gamma=2\) into the \(\mathrm{L}\)-equation:
\begin{align*}
-3(1)+2\beta+2&=1 \\
-1+2\beta&=1 \\
2\beta&=2 \\
\beta&=1
\end{align*}

Hence
\[
\alpha=1,\qquad \beta=1,\qquad \gamma=2
\]

So the model is
\[
F=\frac12\rho Av^2
\]
\end{enumerate}
\end{djsSolution}


\newpage

% ---- Question 6 ----
\questionitem
% [Source: _QBT___Solns__Dimensional_Analysis, Q6]

A motor pulls a crate of mass $m$ along a rough horizontal floor. The equation below is suggested to model the pulling force exerted by the motor, $F$, while the crate moves in a straight line.
\[
F = k_1 m v\dfrac{\mathrm{d}v}{\mathrm{d}x} + k_2 \mu mg + k_3 \dfrac{Q}{x}
\]
in which \\

$\bullet$ $v$ is the speed of the crate when it has travelled distance $x$

$\bullet$ $\mu$ is the coefficient of friction between the crate and the floor

$\bullet$ $g$ is the acceleration due to gravity

$\bullet$ $Q$ is the total thermal energy transferred to the surroundings after the crate has moved distance $x$

$\bullet$ $k_1$, $k_2$ and $k_3$ are dimensionless constants \\

Determine whether the equation is dimensionally consistent. \marks{6} \\

\vspace*{10pt}

\begin{djsSolution}
For the equation to be dimensionally consistent, every term must have the dimensions of force:
\[
[F]=\mathrm{MLT^{-2}}
\]

Also, \(k_1\), \(k_2\) and \(k_3\) are dimensionless, and \(\mu\) is dimensionless.

Now consider each term on the right-hand side:
\[
[v]=\mathrm{LT^{-1}}, \qquad [x]=\mathrm{L}
\]

so
\[
\left[\frac{\mathrm{d}v}{\mathrm{d}x}\right]=\frac{[v]}{[x]}=\frac{\mathrm{LT^{-1}}}{\mathrm{L}}=\mathrm{T^{-1}}
\]

For the first term:
\begin{align*}
\left[k_1mv\frac{\mathrm{d}v}{\mathrm{d}x}\right]
&=[m][v]\left[\frac{\mathrm{d}v}{\mathrm{d}x}\right] \\
&=\mathrm{M}\times\mathrm{LT^{-1}}\times\mathrm{T^{-1}} \\
&=\mathrm{MLT^{-2}}
\end{align*}

For the second term:
\begin{align*}
\left[k_2\mu mg\right]
&=[m][g] \\
&=\mathrm{M}\times\mathrm{LT^{-2}} \\
&=\mathrm{MLT^{-2}}
\end{align*}

For the third term, \(Q\) is thermal energy transferred, so it has the dimensions of energy:
\[
[Q]=\mathrm{ML^{2}T^{-2}}
\]

Therefore
\begin{align*}
\left[k_3\frac{Q}{x}\right]
&=\frac{[Q]}{[x]} \\
&=\frac{\mathrm{ML^{2}T^{-2}}}{\mathrm{L}} \\
&=\mathrm{MLT^{-2}}
\end{align*}

Each term on the right-hand side has dimensions \(\mathrm{MLT^{-2}}\), the same as force.

Therefore the equation is dimensionally consistent.
\end{djsSolution}


\newpage

% ---- Question 7 ----
\questionitem
% [Source: _QBT___Solns__Dimensional_Analysis, Q7]

The speed of a particle falling vertically through a resistive medium is modelled by the formula
\[
v=\dfrac{mg}{k}\left(1-e^{-kt/m}\right)
\] 
where \\

$v$ is the speed of the particle

$m$ is the mass of the particle

$g$ is the acceleration due to gravity

$t$ is the time since the particle was released

$k$ is a constant \\

Use dimensional analysis to determine the dimensions of $k$. \marks{2} \\

\vspace*{10pt}

\begin{djsSolution}
For
\[
v=\frac{mg}{k}\left(1-e^{-kt/m}\right)
\]

the factor \(1-e^{-kt/m}\) must be dimensionless.

In particular, the exponent of an exponential must have no dimensions, so
\[
\left[\frac{kt}{m}\right]=1
\]

Using
\[
[t]=\mathrm{T},\qquad [m]=\mathrm{M}
\]

we get
\[
\frac{[k][t]}{[m]}=1
\]

so
\[
\frac{[k]\mathrm{T}}{\mathrm{M}}=1
\]

Hence
\[
[k]=\mathrm{MT^{-1}}
\]

so the SI unit of \(k\) is \(\mathrm{kg\,s^{-1}}\).
\end{djsSolution}


\newpage

% ---- Question 8 ----
\questionitem
% [Source: _QBT___Solns__Dimensional_Analysis, Q8]

A comet of mass $m$ measured in $\mathrm{kg}$ follows an elliptical orbit around a planet. The semi-major axis of the orbit is $a$ measured in metres and the orbital period is $T$ measured in seconds. The eccentricity $e$ of the orbit is dimensionless. \\

A constant $\mu$ associated with the planet is defined by
\[
T^2 = \frac{4\pi^2 a^3}{\mu}
\]
and the angular momentum of the comet is given by
\[
H = m\sqrt{\mu a(1-e^2)}
\]
By using dimensional analysis, find the dimensions of:

\begin{enumerate}[label=\textbf{(\alph*)}]
  \item $\mu$ \marks{2}
  \item $H$ \marks{2}
\end{enumerate}

\vspace*{10pt}

\begin{djsSolution}
\begin{enumerate}[label=\textbf{(\alph*)}, start=1]
  \item In
  \[
  T^2=\frac{4\pi^2a^3}{\mu}
  \]

  the factor \(4\pi^2\) is dimensionless, so
  \[
  [T]^2=\frac{[a]^3}{[\mu]}
  \]

  Using \([T]=\mathrm{T}\) and \([a]=\mathrm{L}\) gives
  \[
  \mathrm{T^2}=\frac{\mathrm{L^3}}{[\mu]}
  \]

  so
  \[
  [\mu]=\frac{\mathrm{L^3}}{\mathrm{T^2}}=\mathrm{L^3T^{-2}}
  \]

  \item The angular momentum is
  \[
  H=m\sqrt{\mu a(1-e^2)}
  \]

  Here \([m]=\mathrm{M}\), and since \(e\) is dimensionless, \((1-e^2)\) is also dimensionless.

  Therefore
  \[
  [H]=[m]\sqrt{[\mu][a]}
  \]

  Using \([\mu]=\mathrm{L^3T^{-2}}\) from part (a) and \([a]=\mathrm{L}\) gives
  \begin{align*}
  [H] &= \mathrm{M}\sqrt{(\mathrm{L^3T^{-2}})(\mathrm{L})} \\
       &= \mathrm{M}\sqrt{\mathrm{L^4T^{-2}}} \\
       &= \mathrm{ML^2T^{-1}}
  \end{align*}
\end{enumerate}
\end{djsSolution}


\newpage

% ---- Question 9 ----
\questionitem
% [Source: _QBT___Solns__Dimensional_Analysis, Q9]

A small ball of mass $m\,\mathrm{kg}$ is projected vertically upwards with speed $u\,\mathrm{m\,s^{-1}}$. During the motion, when its speed is $v\,\mathrm{m\,s^{-1}}$, it experiences an air resistance of magnitude $kv^2\,\mathrm{N}$, where $k$ is a constant. \\
\begin{enumerate}[label=\textbf{(\alph*)}]
  \item Determine the dimensions of $k$. \marks{2} \\
\end{enumerate}
The greatest height reached by the ball is $H\,\mathrm{m}$. A formula approximating $H$ is
\[
H = \dfrac{u^2}{2g} - \dfrac{ku^4}{4mg^2} + \dfrac{1}{6}k^2 m^\alpha u^\beta g^\gamma
\]
where $g\,\mathrm{m\,s^{-2}}$ is the acceleration due to gravity. \\
\begin{enumerate}[label=\textbf{(\alph*)}, resume]
  \item Use dimensional analysis to find $\alpha$, $\beta$ and $\gamma$. \marks{4} \\
\end{enumerate}

\vspace*{10pt}

\begin{djsSolution}
\begin{enumerate}[label=\textbf{(\alph*)}, start=1]
  \item The air resistance has magnitude \(kv^2\), and this is a force, so
\[
[k][v]^2=[F]
\]

Now
\[
[F]=\mathrm{MLT^{-2}}, \qquad [v]=\mathrm{LT^{-1}}
\]

So
\begin{align*}
[k] &= \frac{[F]}{[v]^2}\\
&= \frac{\mathrm{MLT^{-2}}}{(\mathrm{LT^{-1}})^2}\\
&= \frac{\mathrm{MLT^{-2}}}{\mathrm{L^2T^{-2}}}\\
&= \mathrm{ML^{-1}}
\end{align*}


Hence the SI unit of \(k\) is \(\mathrm{kg\,m^{-1}}\).

  \item Since \(H\) is a height,
\[
[H]=\mathrm{L}
\]

Therefore each term in
\[
H = \dfrac{u^2}{2g} - \dfrac{ku^4}{4mg^2} + \dfrac{1}{6}k^2 m^\alpha u^\beta g^\gamma
\]

must have dimensions of length.

So the third term must satisfy
\[
[k^2m^\alpha u^\beta g^\gamma]=\mathrm{L}
\]

Using
\[
[k]=\mathrm{ML^{-1}}, \qquad [m]=\mathrm{M}, \qquad [u]=\mathrm{LT^{-1}}, \qquad [g]=\mathrm{LT^{-2}}
\]

we get
\begin{align*}
[k^2m^\alpha u^\beta g^\gamma]
&=(\mathrm{ML^{-1}})^2(\mathrm{M})^\alpha(\mathrm{LT^{-1}})^\beta(\mathrm{LT^{-2}})^\gamma \\
&= \mathrm{M}^{2+\alpha}\mathrm{L}^{-2+\beta+\gamma}\mathrm{T}^{-\beta-2\gamma}
\end{align*}

This must equal \(\mathrm{M^0L^1T^0}\), so equating powers gives
\begin{align*}
2+\alpha &= 0 \\
-2+\beta+\gamma &= 1 \\
-\beta-2\gamma &= 0
\end{align*}

The first equation gives
\[
\alpha=-2
\]

The third equation gives
\[
\beta=-2\gamma
\]

Substitute into the second equation:
\begin{align*}
-2+(-2\gamma)+\gamma &= 1 \\
-2-\gamma &= 1 \\
-\gamma &= 3 \\
\gamma &= -3
\end{align*}

Then
\[
\beta=-2(-3)=6
\]

Hence
\[
\alpha=-2, \qquad \beta=6, \qquad \gamma=-3
\]
\end{enumerate}
\end{djsSolution}

\newpage 

% ---- Question 10 ----
\questionitem
% [Source: _QBT___Solns__Dimensional_Analysis, Q10]

\begin{figure}[H]
    \centering
\begin{tikzpicture}[scale=1.1]
\def\Height{3.0}
\def\Range{8.2}
\def\Angle{38}
\def\VelLen{2.3}
\def\DimDrop{0.8}
\def\DimOff{0.55}
\pgfmathsetmacro{\TanA}{tan(\Angle)}
\pgfmathsetmacro{\Coeff}{(\Height+\TanA*\Range)/(\Range*\Range)}

\coordinate (B) at (0,0);
\coordinate (O) at (0,\Height);
\coordinate (A) at (\Range,0);
\coordinate (Href) at ($(O)+(1.6,0)$);
\coordinate (V) at ($(O)+(\Angle:\VelLen)$);
\coordinate (R1) at ($(B)+(0,-\DimDrop)$);
\coordinate (R2) at ($(A)+(0,-\DimDrop)$);
\coordinate (h1) at ($(B)+(-\DimOff,0)$);
\coordinate (h2) at ($(O)+(-\DimOff,0)$);
\coordinate (h3) at ($(O)+(2*\DimOff,0)$);

\draw[thick] (-1.2,0) -- (\Range+1.2,0);
\draw[thick] (B) -- (O);
\draw[thick, black, domain=0:\Range, samples=200]
  plot (\x, {\Height+\TanA*\x-\Coeff*\x*\x});

\draw[-{Stealth[length=3mm]}, thick] (O) -- (V)
  node[pos=0.65, above] {$u$};

\pic["$\theta$", draw, angle radius=7mm, angle eccentricity=1.35]
  {angle=Href--O--V};

\draw[dashed] (B) -- (h1);
\draw[dashed] (O) -- (h3);
\draw[{Stealth[length=3mm]}-{Stealth[length=3mm]}, thick] (h1) -- (h2)
  node[midway, left] {$h$};

\draw[{Stealth[length=3mm]}-{Stealth[length=3mm]}, thick] (R1) -- (R2)
  node[midway, below] {$R$};

\fill[black] (A) circle (1.2pt);

\node[above left] at (O) {$O$};
\node[below left] at (B) {$B$};
\node[below right] at (A) {$A$};
\end{tikzpicture}
\end{figure}


A particle is projected with speed $u$ from a point $O$ at the top of a cliff of height $h$ above horizontal ground. The angle of projection is $\theta$ above the horizontal. The particle lands at the point $A$ on the ground, and $B$ is the foot of the cliff, so that $BA = R$, as shown in the diagram. \\

You are given that
\[
R = \frac{u^\alpha \cos\theta}{g^\beta}\left(u^\gamma\sin\theta + \sqrt{u^\delta\sin^2\theta + 2gh}\right)
\]
Use dimensional analysis to find $\alpha$, $\beta$, $\gamma$ and $\delta$. \marks{4} \\

\vspace*{10pt}

\begin{djsSolution}
The dimensions of the quantities are
\[
[R]=\mathrm{L},\qquad [u]=\mathrm{LT^{-1}},\qquad [g]=\mathrm{LT^{-2}},\qquad [h]=\mathrm{L}
\]

and \(\sin\theta\) and \(\cos\theta\) are dimensionless.

First look at the expression inside the square root. Terms being added must have the same dimensions, so
\[
[u^\delta\sin^2\theta]=[2gh]
\]

Now
\[
[2gh]=[g][h]=(\mathrm{LT^{-2}})(\mathrm{L})=\mathrm{L^2T^{-2}}
\]

So
\[
(\mathrm{LT^{-1}})^\delta=\mathrm{L^2T^{-2}}
\]

Equating powers of \(\mathrm{L}\) and \(\mathrm{T}\) gives
\[
\delta=2
\]

Hence the whole square-root term has dimensions
\[
\sqrt{\mathrm{L^2T^{-2}}}=\mathrm{LT^{-1}}
\]

Therefore \(u^\gamma\sin\theta\) must also have dimensions \(\mathrm{LT^{-1}}\), so
\[
(\mathrm{LT^{-1}})^\gamma=\mathrm{LT^{-1}}
\]

which gives
\[
\gamma=1
\]

Now the bracket
\[
u^\gamma\sin\theta+\sqrt{u^\delta\sin^2\theta+2gh}
\]

has dimensions \(\mathrm{LT^{-1}}\). So
\begin{align*}
[R]
&=\left[\frac{u^\alpha\cos\theta}{g^\beta}\right]\left[ u^\gamma\sin\theta+\sqrt{u^\delta\sin^2\theta+2gh}\right] \\
\mathrm{L}
&=\frac{(\mathrm{LT^{-1}})^\alpha}{(\mathrm{LT^{-2}})^\beta}(\mathrm{LT^{-1}}) \\
&=\mathrm{L}^{\alpha-\beta+1}\mathrm{T}^{-\alpha+2\beta-1}
\end{align*}

Equating powers of \(\mathrm{L}\) and \(\mathrm{T}\) with \(\mathrm{L^1T^0}\):
\[
\alpha-\beta+1=1,\qquad -\alpha+2\beta-1=0
\]

The first equation gives
\[
\alpha=\beta
\]

Substitute into the second:
\[
-\beta+2\beta-1=0
\]

so
\[
\beta=1
\]

and therefore
\[
\alpha=1
\]

So the values are
\[
\alpha=1,\qquad \beta=1,\qquad \gamma=1,\qquad \delta=2
\]
\end{djsSolution}


\end{enumerate}
\end{document}
